%\documentclass[hyperref={pdfpagelabels=false},slidetop,9pt]{beamer}
\documentclass[slidetop,8pt,aspectratio=169]{beamer}
%\usepackage{mathpazo}            % Mathpazo  symbols
%\renewcommand {\rmdefault}{ibh} % Bradhley  Hand  ITC
\usepackage{arev}
\usepackage[T1]{fontenc}
\usepackage[utf8]{inputenc}
\newcommand{\id}{71}
\newcommand{\nom}{Théorie des mécanismes}
\newcommand{\sequence}{04}
\newcommand{\nomsequence}{Liaisons entre les solides}
\newcommand{\num}{02}
\newcommand{\type}{KH}
\newcommand{\descrip}{Liaisons équivalentes, hyperstatisme, liaisons en série et en parallèle, théorie des graphes}
\newcommand{\competences}{B2-12: Proposer une modélisation des liaisons avec leurs caractéristiques géométriques. \\ &  B2-13: Proposer un modèle cinématique paramétré à partir d'un système réel, d'une maquette numérique ou d'u \\ &  B2-17: Simplifier un modèle de mécanisme. \\ &  B2-18: Modifier un modèle pour le rendre isostatique. \\ &  C1-04: Proposer une démarche permettant d'obtenir une loi entrée-sortie géométrique.  \\ &  C2-05: Caractériser le mouvement d'un repère par rapport à un autre repère. \\ &  C2-06: Déterminer les relations entre les grandeurs géométriques ou cinématiques. }
\newcommand{\nbcomp}{7}
\newcommand{\systemes}{}
\newcommand{\systemesnum}{}
\newcommand{\systemessansaccent}{}
\newcommand{\ilot}{2}
\newcommand{\ilotstr}{02}
\newcommand{\dossierilot}{\detokenize{Ilot_02 }}

%\usepackage{etex}
\usepackage{beaudiomer}
\usepackage{tikz}
\usepackage[french]{babel}
\usetikzlibrary{automata,positioning,fit,babel,tikzmark}
\usepackage{animate}
\usepackage[european, cuteinductors]{circuitikz}
\usepackage{pgf}
\usepackage[all]{xy}
\usepackage{pgfpages}
\usepackage{graphbox}
\usepackage{pdfpages}
%\usepackage[adobe-utopia]{mathdesign}
\usepackage{ifthen}
\usepackage{cancel}
\usepackage{framed}
\usepackage{subfig}
\usepackage{tabularx}
\usepackage{setspace}
%\usepackage{soul}
\usepackage{schemabloc}
\usepackage{eqnarray}
\usepackage[dot, phantomtext]{dashundergaps}
\usepackage{media9}
\usepackage{multimedia}
\usepackage{textcomp}
%\usepackage{txfonts}
%\usepackage[upright]{txgreeks}
%\renewcommand\familydefault{fjd} % emerald package
%\usepackage{mathastext}

\author{Renaud Costadoat}
\institute{Lycée Dorian}

\usepackage{multido}
\usepackage{multirow}
\usepackage{multicol} % Portions de texte en colonnes
\usepackage{flafter}%floatants après la référence

\usepackage{color}
\usepackage{xcolor}
\usepackage{tcolorbox}
\usepackage{colortbl}

\usepackage[gen]{eurosym}
\usepackage{tikz}
%\usepackage{pstricks,pst-node,pst-tree,pst-solides3d}
%\usepackage{lmodern}
\usetheme{renaudvideo}
\usepackage{times}
%\usepackage{amsmath}
\usepackage{verbatim}
%\usepackage{moreverb}
\usetikzlibrary{shapes}
\usepackage{graphicx}
\usepackage{psfrag}
\usepackage{wrapfig}
\usepackage{etoolbox}

\definecolor{gris25}{gray}{0.75}
\definecolor{bleu}{RGB}{18,33,98}
\definecolor{bleuf}{RGB}{42,94,171}
\definecolor{bleuc}{RGB}{231,239,247}
\definecolor{rougef}{RGB}{185,18,27}
\definecolor{rougec}{RGB}{255,188,204}%255,230,231
\definecolor{vertf}{RGB}{21,151,19}
\definecolor{vertc}{RGB}{220,255,191}

\addtobeamertemplate{frametitle}{\hspace*{1em}}{}\setlength\parindent{24pt}
\parskip 7.2pt
\parindent 8pt

\newenvironment{rem}[1][\hsize]%
{%
    \def\FrameCommand
   {%
\rotatebox{90}{\textit{\textsf{Remarque}}} 
       {\color{bleuf}\vrule width 3pt}%
       \hspace{0pt}%must no space.
       \fboxsep=\FrameSep\colorbox{bleuc}%
  }%
    \MakeFramed{\hsize#1\advance\hsize-\width\FrameRestore}%
}%
{\endMakeFramed}%


\newenvironment{savoir}[1][\hsize]%
{%
    \def\FrameCommand
    {%
\rotatebox{90}{\textit{\textsf{Savoir}}} 
        {\color{bleuf}\vrule width 3pt}%
        \hspace{0pt}%must no space.
        \fboxsep=\FrameSep\colorbox{bleuc}%
    }%
    \MakeFramed{\hsize#1\advance\hsize-\width\FrameRestore}%
}%
{\endMakeFramed}%

\newenvironment{prob}[1][\hsize]%
{%
    \def\FrameCommand%
    {%
\rotatebox{90}{\textit{\textsf{Problematique}}} 
        {\color{rougef}\vrule width 3pt}%
        \hspace{0pt}%must no space.
        \fboxsep=\FrameSep\colorbox{rougec}%
    }%
    \MakeFramed{\hsize#1\advance\hsize-\width\FrameRestore}%
}%
{\endMakeFramed}%

\newenvironment{obj}[1][\hsize]%
{%
    \def\FrameCommand%
    {%
\rotatebox{90}{\textit{\textsf{Objectif}}} 
        {\color{vertf}\vrule width 3pt}%
        \hspace{0pt}%must no space.
        \fboxsep=\FrameSep\colorbox{vertc}%
    }%
    \MakeFramed{\hsize#1\advance\hsize-\width\FrameRestore}%
}%
{\endMakeFramed}%

\newenvironment{defi}[1][\hsize]%
{%
    \def\FrameCommand%
    {%
\rotatebox{90}{\textit{\textsf{Definition}}} 
        {\color{bleuf}\vrule width 3pt}%
        \hspace{0pt}%must no space.
        \fboxsep=\FrameSep\colorbox{rougec}%
    }%
    \MakeFramed{\hsize#1\advance\hsize-\width\FrameRestore}%
}%
{\endMakeFramed}%


\newenvironment{hypo}[1][\hsize]%
{%
    \def\FrameCommand%
    {%
\rotatebox{90}{\textit{\textsf{Hypothèse\\}}} 
        {\color{bleuf}\vrule width 3pt}%
        \hspace{0pt}%must no space.
        \fboxsep=\FrameSep\colorbox{bleuc}%
    }%
    \MakeFramed{\hsize#1\advance\hsize-\width\FrameRestore}%
}%
{\endMakeFramed}%


\newenvironment{prop}[1][\hsize]%
{%
    \def\FrameCommand%
    {%
\rotatebox{90}{\textit{\textsf{Propriété}}} 
        {\color{bleuf}\vrule width 3pt}%
        \hspace{0pt}%must no space.
        \fboxsep=\FrameSep\colorbox{bleuc}%
    }%
    \MakeFramed{\hsize#1\advance\hsize-\width\FrameRestore}%
}%
{\endMakeFramed}%

\newenvironment{props}[1][\hsize]%
{%
    \def\FrameCommand%
    {%
\rotatebox{90}{\textit{\textsf{Propriétés}}} 
        {\color{bleuf}\vrule width 3pt}%
        \hspace{0pt}%must no space.
        \fboxsep=\FrameSep\colorbox{bleuc}%
    }%
    \MakeFramed{\hsize#1\advance\hsize-\width\FrameRestore}%
}%
{\endMakeFramed}%

\newenvironment{exemple}[1][\hsize]%
{%
    \def\FrameCommand%
    {%
\rotatebox{90}{\textit{\textsf{Exemple}}} 
        {\color{vertf}\vrule width 3pt}%
        \hspace{0pt}%must no space.
        \fboxsep=\FrameSep\colorbox{vertc}%
    }%
    \MakeFramed{\hsize#1\advance\hsize-\width\FrameRestore}%
}%
{\endMakeFramed}%

\newenvironment{resultat}[1][\hsize]%
{%
    \def\FrameCommand%
    {%
\rotatebox{90}{\textit{\textsf{Résultat}}} 
        {\color{rougef}\vrule width 3pt}%
%        {\color{bleuf}\vrule width 3pt}%
        \hspace{0pt}%must no space.
        \fboxsep=\FrameSep\colorbox{rougec}%
    }%
    \MakeFramed{\hsize#1\advance\hsize-\width\FrameRestore}%
}%
{\endMakeFramed}%

\newenvironment{methode}[1][\hsize]%
{%
    \def\FrameCommand%
    {%
\rotatebox{90}{\textit{\textsf{Méthode\\}}} 
        {\color{rougef}\vrule width 3pt}%
        \hspace{0pt}%must no space.
        \fboxsep=\FrameSep\colorbox{rougec}%
    }%
    \MakeFramed{\hsize#1\advance\hsize-\width\FrameRestore}%
}%
{\endMakeFramed}%

\newenvironment{theo}[1][\hsize]%
{%
    \def\FrameCommand%
    {%
\rotatebox{90}{\textit{\textsf{Théorème\\}}} 
        {\color{rougef}\vrule width 3pt}%
        \hspace{0pt}%must no space.
        \fboxsep=\FrameSep\colorbox{rougec}%
    }%
    \MakeFramed{\hsize#1\advance\hsize-\width\FrameRestore}%
}%
{\endMakeFramed}%

\newenvironment{warn}[1][\hsize]%
{%
    \def\FrameCommand%
    {%
\rotatebox{90}{\textit{\textsf{Attention\\}}} 
        {\color{rougef}\vrule width 3pt}%
        \hspace{0pt}%must no space.
        \fboxsep=\FrameSep\colorbox{rougec}%
    }%
    \MakeFramed{\hsize#1\advance\hsize-\width\FrameRestore}%
}%
{\endMakeFramed}%


\newcommand{\feuilleDR}[1]{
	\begin{tikzpicture}
		\draw[gray!30](0,0)grid[step=0.5cm](\linewidth,#1);
	\end{tikzpicture}
}

% \usepackage{pstricks}
%\usepackage{minitoc}
% \setcounter{minitocdepth}{4}

\setcounter{tocdepth}{2}

% \mtcselectlanguage{french} 

%\usepackage{draftcopy}% "Brouillon"
% \usepackage{floatflt}
\usepackage{psfrag}
%\usepackage{listings} % Permet d'insérer du code de programmation
\renewcommand{\baselinestretch}{1.2}

% Changer la num�rotation des figures :
% ------------------------------------
% \makeatletter
% \renewcommand{\thefigure}{\ifnum \c@section>\z@ \thesection.\fi
%  \@arabic\c@figure}
% \@addtoreset{figure}{section}
% \makeatother
 


%%%%%%%%%%%%
% Définition des vecteurs %
%%%%%%%%%%%%
 \newcommand{\vect}[1]{\overrightarrow{#1}}

%%%%%%%%%%%%
% Définition des torseusr %
%%%%%%%%%%%%

 \newcommand{\torseur}[1]{%
\left\{{#1}\right\}
}

\newcommand{\torseurcin}[3]{%
\left\{\mathcal{#1} \left(#2/#3 \right) \right\}
}

\newcommand{\torseurstat}[3]{%
\left\{\mathcal{#1} \left(#2\rightarrow #3 \right) \right\}
}

 \newcommand{\torseurc}[8]{%
%\left\{#1 \right\}=
\left\{
{#1}
\right\}
 = 
\left\{%
\begin{array}{cc}%
{#2} & {#5}\\%
{#3} & {#6}\\%
{#4} & {#7}\\%
\end{array}%
\right\}_{#8}%
}

 \newcommand{\torseurcol}[7]{
\left\{%
\begin{array}{cc}%
{#1} & {#4}\\%
{#2} & {#5}\\%
{#3} & {#6}\\%
\end{array}%
\right\}_{#7}%
}

 \newcommand{\torseurl}[3]{%
%\left\{\mathcal{#1}\right\}_{#2}=%
\left\{%
\begin{array}{l}%
{#1} \\%
{#2} %
\end{array}%
\right\}_{#3}%
}

 \newcommand{\vectv}[3]{%
\vect{V\left( {#1} \in {#2}/{#3}\right)}
}


\newcommand{\vectf}[2]{%
\vect{R\left( {#1} \rightarrow {#2}\right)}
}

\newcommand{\vectm}[3]{%
\vect{\mathcal{M}\left( {#1}, {#2} \rightarrow {#3}\right)}
}


 \newcommand{\vectg}[3]{%
\vect{\Gamma \left( {#1} \in {#2}/{#3}\right)}
}

 \newcommand{\vecto}[2]{%
\vect{\Omega\left( {#1}/{#2}\right)}
}

\newcommand{\reponse}[1][4]
{
\multido{}{#1}
{
\begin{center}
\makebox[0.9\linewidth]{\dotfill} \end{center}
}}


% }$$\left\{\mathcal{#1} \right\}_{#2} =%
% \left\{%
% \begin{array}{c}%
%  #3 \\%
%  #4 %
% \end{array}%
% \right\}_{#5}}


%  ------------------------------------------
% | Modification du formatage des sections : | 
%  ------------------------------------------

% Grands titres :
% ---------------

\newcommand{\titre}[1]{%
\begin{center}
      \bigskip
      \rule{\textwidth}{1pt}
      \par\vspace{0.1cm}
      
      \textbf{\large #1}
      \par\rule{\textwidth}{1pt}
    \end{center}
    \bigskip
  }

% Supprime le numéro du chapitre dans la numérotation des sections:
% -----------------------------------------------------------------
\makeatletter
\renewcommand{\thesection}{\@arabic\c@section}
\makeatother

\newtcbox{\mybox}{on line,
  colframe=rougef,colback=rougec,
  boxrule=0.5pt,arc=4pt,boxsep=0pt,left=3pt,right=3pt,top=3pt,bottom=3pt}
 
\newtcbox{\myboxt}{on line,
  colframe=white!0,colback=white!0,opacityframe=0, opacityback=0,
  boxrule=0.5pt,arc=4pt,boxsep=0pt,left=3pt,right=3pt,top=3pt,bottom=3pt,
  }
 



% \titleformat{\chapter}[display]
% {\normalfont\Large\filcenter}
% {}
% {1pc}
% {\titlerule[1pt]
%   \vspace{1pc}%
%   \Huge}[\vspace{1ex}%
% \titlerule]


%%%% Chapitres Comme PY Pechard %%%%%%%%%
% numéro du chapitre
\DeclareFixedFont{\chapnumfont}{OT1}{phv}{b}{n}{80pt}
% pour le mot " Chapitre "
\DeclareFixedFont{\chapchapfont}{OT1}{phv}{m}{it}{40pt}
% pour le titre
\DeclareFixedFont{\chaptitfont}{T1}{phv}{b}{n}{25pt}

\definecolor{gris}{gray}{0.75}
\setbeamertemplate{section in toc}[sections numbered]

\newlength{\RoundedBoxWidth}
\newsavebox{\GrayRoundedBox}
\newenvironment{GrayBox}[1][\dimexpr\textwidth-4.5ex]%
   {\setlength{\RoundedBoxWidth}{\dimexpr#1}
    \begin{lrbox}{\GrayRoundedBox}
       \begin{minipage}{\RoundedBoxWidth}}%
   {   \end{minipage}
    \end{lrbox}
    \begin{center}
    \begin{tikzpicture}%
       \draw node[draw=bleuf,fill=bleuc,rounded corners,%
             inner sep=2ex,text width=\RoundedBoxWidth]%
             {\usebox{\GrayRoundedBox}};
    \end{tikzpicture}
    \end{center}}
    
\ifdef{\prive}{\pgfpagesuselayout{2 on 1}[a4paper,border shrink=0mm]}
\ifdef{\prive}{\setbeamertemplate{navigation symbols}{}}
\setbeamertemplate{itemize item}[ball]
%\setbeamertemplate{blocks}[rounded]%[shadow=true]
\setbeamercolor{block title}{fg=white,bg=grisf}        % titre block normal 
\setbeamercolor{block body}{fg=grisf,bg=grisc!50}      % corps block normal
\setbeamercolor{block body alerted}{fg=white,bg=warning}   % idem pour un block alerte

\date{S\sequence \ - \type\num}

\begin{document}%\mathversion{bold}
\shorthandoff{:!}
\bibliographystyle{abbrvnat-fr}

\setbeamertemplate{footline}[pagetitre]


\title{Fonctions de transfert dans le domaine de Laplace}

\setbeamertemplate{headline}[pagetitre]
\frame{\titlepage}

\usebackgroundtemplate%
{%
    \centering\includegraphics[width=\paperwidth]{../../headers/tableau}%
}

\setbeamertemplate{headline}[pagetitre]
\frame{
\frametitle{\inserttitle}

\begin{obj}
Déterminer la fonction de transfert qui régit le comportement d'un S.L.C.I.
\end{obj}

\begin{minipage}{0.46\linewidth}
\tableofcontents
\end{minipage}\hfill
\begin{minipage}{0.46\linewidth}
\begin{figure}[ht!]
\centering\includegraphics[width=0.5\linewidth]<2>{img/moteur}
\end{figure}
\end{minipage}
}

\section[Mise en équation]{Mise en équation (Moteur à Courant Continu $\frac{\omega_m}{u_m}$)}

\setbeamertemplate{headline}[renaud theme]
{\frame{
\frametitle{Mise en équation}

\begin{minipage}{0.47\linewidth}
\begin{eqnarray}
	\color<1>{blue}{u_m(t)=R_m\cdot i_m(t)+L_m\cdot \frac{d\ i_m(t)}{dt}+e(t)} \\
	\color<2>{blue}{e(t)=K_e\cdot \omega_m(t)} \\
	\color<2>{blue}{c_m(t)=K_i\cdot i_m(t)} \\
	\color<3>{blue}{J\cdot \frac{d\omega_m(t)}{dt}=c_m(t)}
\end{eqnarray}
\end{minipage}\hfill
\begin{minipage}{0.47\linewidth}
 \begin{overlayarea}{6.5cm}{4cm}
\only<1>{
 \begin{circuitikz}
\draw (0,3) to [R, l=$R$, i=$i_m$] (2.5,3);
\draw (2.5,3) to [L, l=$L$] (5,3);
\draw (5,0) to [V, v=$e(t)$] (5,3);
\draw (5,0) -- (0,0);
\draw[-triangle 45] (0,0.5) --  node[left] {$u_m(t)$} (0,2.5) ;
\end{circuitikz}}
\only<3>{ \centering\begin{tikzpicture}[scale=0.8]
\draw[fill=blue!50](0,0) circle (1 and 2);
\draw[top color=blue!25,bottom color=black,middle color=blue!50] (-0.5,2) arc (90:270:1 and 2) -- ++(0.5,0) arc (-90:-270:1 and 2) -- cycle;
\draw[top color=white,bottom color=black!70] (0,3mm) arc (90:270:1.5mm and 3mm)--++(3cm,0) arc (-90:-270:1.5mm and 3mm)-- cycle;
\draw (3cm,3mm) arc (90:-90:1.5mm and 3mm);
\end{tikzpicture}}
\end{overlayarea}
\end{minipage}
}}

\setbeamertemplate{headline}[renaud theme]
{\frame{
\frametitle{Équation différentielle}
\begin{eqnarray}
\frac{u_m(t)}{K_e}=\frac{L_m\cdot J}{K_i\cdot K_e}\cdot\frac{d^2\omega_m(t)}{dt^2}+\frac{R_m\cdot J}{K_i\cdot K_e}\cdot\frac{d\omega_m(t)}{dt}+\omega_m(t)
\end{eqnarray}

Équation différentielle d'\textbf{ordre 2 avec second membre}. Méthode de résolution:
\begin{itemize}
 \item Recherche de la solution de l'équation sans second membre à l'aide du \textbf{discriminant},
 \item Recherche d'une \textbf{solution particulière} afin d'obtenir la solution générale.
\end{itemize}

Autre solution possible: \textbf{La transformée de Laplace} qui permet de résoudre des équations différentielles \textbf{linéaires} à \textbf{coefficients constants}.


}}

\section{Transformées de Laplace}

\setbeamertemplate{headline}[renaud theme]
{\frame{
\frametitle{Passage dans le domaine de Laplace}

\setcounter{equation}{0}

\begin{minipage}{0.47\linewidth}
\begin{eqnarray}
	u_m(t)=R_m\cdot i_m(t)+L_m\cdot \frac{d\ i_m(t)}{dt}+e(t) \\
	e(t)=K_e\cdot \omega_m(t) \\
	c_m(t)=K_i\cdot i_m(t) \\
	J\cdot \frac{d\omega_m(t)}{dt}=c_m(t)
\end{eqnarray}
\end{minipage}\hfill
\begin{minipage}{0.47\linewidth}
\only<2>{
Théorème de la dérivée première:\\ ~\ \\
	$L\left[\frac{df(t)}{dt}\right]=p\cdot F(p)-f(0^+)$
	Hypothèse:
	\begin{itemize}
	 \item Conditions initiales nulles (obligatoire pour la Fonction de Transfert)
	\end{itemize}
}
\only<3>{\begin{eqnarray}
	U_m(p)=R_m\cdot I_m(p)+L_m\cdot p\cdot I_m+E(p) \\
	E(p)=K_e\cdot \Omega_m(p) \\
	C_m(p)=K_i\cdot I_m(p) \\
	J\cdot p\cdot \Omega_m(p)=C_m(p)
\end{eqnarray}}
\end{minipage}
}}

\section{Fonction de transfert}

\setbeamertemplate{headline}[renaud theme]
{\frame{
\frametitle{Fonction de transfert}

\setcounter{equation}{4}

\only<1-3>{\begin{eqnarray}
	{\color<2->{vertf}{U_m(p)}}=R_m\cdot {\color<2->{red}{I_m(p)}}+L_m\cdot p\cdot {\color<2->{red}{I_m(p)}}+{\color<2->{red}{E(p)}} \\
	{\color<2->{red}{E(p)}}=K_e\cdot {\color<2->{vertf}{\Omega_m(p)}} \\
	{\color<2->{red}{C_m(p)}}=K_i\cdot {\color<2->{red}{I_m(p)}} \\
	J\cdot p\cdot {\color<2->{vertf}{\Omega_m(p)}}={\color<2->{red}{C_m(p)}}
\end{eqnarray}}
\only<3-6>{
\begin{eqnarray}
	\only<3>{{\color<2->{red}{I_m(p)}}=\frac{\color<2->{red}{C_m(p)}}{K_i} et\ {\color<2->{red}{E(p)}}=K_e\cdot {\color<2->{vertf}{\Omega_m(p)}} \Rightarrow} {\setcounter{equation}{9}\color<3->{vertf}{U_m(p)}}=R_m\cdot \frac{\color<2->{red}{C_m(p)}}{K_i}+L_m\cdot p\cdot \frac{\color<2->{red}{C_m(p)}}{K_i}+K_e\cdot {\color<2->{vertf}{\Omega_m(p)}}
\end{eqnarray}}
\only<5-7>{
\begin{eqnarray}
	\only<5>{{\color<2->{red}{C_m(p)}}=J\cdot p\cdot {\color<2->{vertf}{\Omega_m(p)}}\ (12) \Rightarrow} {\setcounter{equation}{10}\color<3->{vertf}{U_m(p)}}=R_m\cdot \frac{J\cdot p\cdot {\color<2->{vertf}{\Omega_m(p)}}}{K_i}+L_m\cdot p\cdot \frac{J\cdot p\cdot {\color<2->{vertf}{\Omega_m(p)}}}{K_i}+K_e\cdot {\color<2->{vertf}{\Omega_m(p)}}
\end{eqnarray}}
\only<6-8>{
\begin{eqnarray}
	{\setcounter{equation}{11}\color<3->{vertf}{U_m(p)}}=\left(\frac{R_m\cdot J\cdot p}{K_i}+\frac{L_m\cdot J\cdot p^2}{K_i}+K_e\right) \cdot {\color<2->{vertf}{\Omega_m(p)}}
\end{eqnarray}}
\only<8->{
\begin{eqnarray}
	\setcounter{equation}{12}\frac{\color<2->{vertf}{\Omega_m(p)}}{\color<3->{vertf}{U_m(p)}}=\frac{1}{K_e+\frac{R_m\cdot J}{K_i}\cdot p+\frac{L_m\cdot J}{K_i}\cdot p^2} 
\end{eqnarray}}
\only<9->{
Forme canonique:
\begin{eqnarray}
	\setcounter{equation}{13}\frac{\color<2->{vertf}{\Omega_m(p)}}{\color<3->{vertf}{U_m(p)}}=\frac{\frac{1}{K_e}}{
	 \tikzmark{b1} \only<9>{\myboxt{1}} \only<10->{\mybox{1}} +\frac{R_m\cdot J}{K_i\cdot K_e}\cdot p+\frac{L_m\cdot J}{K_i\cdot K_e}\cdot p^2} 
\end{eqnarray}}
\only<10->{\begin{tikzpicture}[overlay,remember picture]
    \draw[arrows=-latex] 
    ( $ (pic cs:b1) +(15pt,-4.5ex) $ ) -- 
    ( $ (pic cs:b1) +(8pt,-1ex) $ );
    \node[anchor=west]
    at ( $ (pic cs:b1) +(12pt,-6ex) $ )
    {Partie constante unitaire};
\end{tikzpicture}}


}}

\end{document}	

